\documentclass[a4paper,12pt]{article}
\usepackage[a4paper,margin=1.8cm]{geometry}
\usepackage{amsmath}
\usepackage{enumitem}
\pagestyle{empty}
\parindent=0pt

\begin{document}

\begin{center}
  {\LARGE\bfseries Answer Key}\\[4pt]
  {\large Drawing Pie Charts with a Protractor}
\end{center}

\vspace{4pt}

\section*{Chart 1 -- Favourite Pets (class of 36 students)}

Completed sector angles: Dog \(120^\circ\), Cat \(90^\circ\), Fish \(80^\circ\), Rabbit \(50^\circ\), Other \(20^\circ\).

\begin{enumerate}[leftmargin=1.6em, itemsep=2pt]
  \item Completed chart; sectors listed above.
  \item Dog, because \(\frac{12}{36}=\frac{1}{3}\).
  \item Cat: 9 students.
\end{enumerate}

\section*{Chart 2 -- Sports Club Memberships (36 members)}

Completed sector angles: Football \(140^\circ\), Swimming \(90^\circ\), Tennis \(60^\circ\), Basketball \(50^\circ\), Other \(20^\circ\).

\begin{enumerate}[leftmargin=1.6em, itemsep=2pt]
  \item Completed chart; sectors listed above.
  \item Swimming has a \(90^\circ\) sector because \(\frac{9}{36}=\frac{1}{4}\).
  \item Tennis fraction \(=\frac{6}{36}=\frac{1}{6}\).
  \item The three smallest sports are Tennis, Basketball and Other:
  \[
  60^\circ+50^\circ+20^\circ=130^\circ.
  \]
\end{enumerate}

\section*{Chart 3 -- How Students Spend Most of Their Pocket Money (36 students)}

Completed sector angles: Food \& Drink \(120^\circ\), Games \(90^\circ\), Clothes \(60^\circ\), Savings \(60^\circ\), Other \(30^\circ\).

\begin{enumerate}[leftmargin=1.6em, itemsep=2pt]
  \item Completed chart; sectors listed above.
  \item Clothes and Savings both have 6 students, so each has a sector angle of \(60^\circ\).
\end{enumerate}

\section*{Chart 4 -- Colours of Cars in a Car Park (60 cars)}

\begin{enumerate}[leftmargin=1.6em, itemsep=2pt]
  \item Completed table:
  \begin{center}
    \renewcommand{\arraystretch}{1.3}
    \begin{tabular}{lccc}
      \hline
      \textbf{Colour} & \textbf{No.} & \textbf{Fraction} & \textbf{Angle}\\
      \hline
      White  & 20 & \(\frac{1}{3}\)  & \(120^\circ\)\\
      Black  & 15 & \(\frac{1}{4}\)  & \(90^\circ\)\\
      Silver & 12 & \(\frac{1}{5}\)  & \(72^\circ\)\\
      Red    &  9 & \(\frac{3}{20}\) & \(54^\circ\)\\
      Other  &  4 & \(\frac{1}{15}\) & \(24^\circ\)\\
      \hline
      Total  & 60 & 1 & \(360^\circ\)\\
      \hline
    \end{tabular}
  \end{center}
  \item Completed pie chart using the sector angles in the table.
  \item White angle:
  \[
  \frac{20}{60}\times360=\frac{1}{3}\times360=120^\circ.
  \]
  \item Black fraction \(=\frac{15}{60}=\frac{1}{4}\).
  \item Angle sum check:
  \[
  120+90+72+54+24=360^\circ.
  \]
\end{enumerate}

\section*{Chart 5 -- Languages Spoken at Home (survey of 90 students)}

\begin{enumerate}[leftmargin=1.6em, itemsep=2pt]
  \item Completed table:
  \begin{center}
    \renewcommand{\arraystretch}{1.3}
    \begin{tabular}{lccc}
      \hline
      \textbf{Language} & \textbf{No.} & \textbf{Fraction} & \textbf{Angle}\\
      \hline
      English & 45 & \(\frac{1}{2}\)  & \(180^\circ\)\\
      Urdu    & 18 & \(\frac{1}{5}\)  & \(72^\circ\)\\
      Polish  &  9 & \(\frac{1}{10}\) & \(36^\circ\)\\
      Punjabi &  9 & \(\frac{1}{10}\) & \(36^\circ\)\\
      Other   &  9 & \(\frac{1}{10}\) & \(36^\circ\)\\
      \hline
      Total   & 90 & 1 & \(360^\circ\)\\
      \hline
    \end{tabular}
  \end{center}
  \item Completed pie chart using: English \(180^\circ\), Urdu \(72^\circ\), Polish/Punjabi/Other \(36^\circ\) each.
  \item Polish, Punjabi and Other each have 9 speakers, so each sector is
  \[
  \frac{9}{90}\times360=\frac{1}{10}\times360=36^\circ.
  \]
  \item English fraction \(=\frac{45}{90}=\frac{1}{2}\).
  \item \textbf{Challenge:} new total \(=100\), new Polish speakers \(=19\), so new Polish sector
  \[
  \frac{19}{100}\times360=68.4^\circ
  \]
  (accept \(68^\circ\) or \(69^\circ\)).
\end{enumerate}

\end{document}